\documentclass[11pt]{article}
\usepackage[a4paper,margin=1in]{geometry}
\usepackage{amsmath,amssymb}
\usepackage{xcolor}
\usepackage{booktabs}
\usepackage{listings}
\usepackage{parskip}
\usepackage{titlesec}
\usepackage{graphicx}
\usepackage{caption}
\usepackage{subcaption}

\titleformat{\section}{\large\bfseries}{\thesection}{0.6em}{}
\lstset{basicstyle=\small\ttfamily, breaklines=true, frame=single, framesep=4pt}

\title{OU angular velocity --- variable cheat-sheet\\
\large \texttt{generate\_path\_integration\_batch} setup block}
\date{2026-05-13}

\begin{document}
\maketitle

The block in lines 292--300 sets up the discretised
Ornstein--Uhlenbeck (OU) process for angular velocity. The
continuous-time SDE is
\begin{equation}
\frac{d\omega}{dt} \;=\; -\alpha\,\omega(t) \;+\; \sigma\,\xi(t),
\qquad \xi(t)\sim\mathrm{white\ noise},
\label{eq:ou}
\end{equation}
with stationary distribution
$\omega\sim\mathcal{N}\bigl(0,\;\sigma^{2}/(2\alpha)\bigr)$ so the
stationary standard deviation is
$\sigma_{\omega}=\sigma/\sqrt{2\alpha}$. The Euler--Maruyama
discretisation that the next loop uses is
\begin{equation}
\omega[t+1] \;=\; \omega[t]
    \;-\;\alpha\,\omega[t]\,\Delta t
    \;+\;\sigma\sqrt{\Delta t}\;\eta[t],
\qquad \eta[t]\sim\mathcal{N}(0,1).
\label{eq:ou-discrete}
\end{equation}

\section*{Variables in the snippet}

\begin{center}\small
\begin{tabular}{lll}
\toprule
\textbf{Code} & \textbf{Math} & \textbf{Meaning / Hulse default} \\
\midrule
\texttt{B}             & $B$                       & Batch size (number of trials per gradient step, default 100). \\
\texttt{T}             & $T$                       & Time steps per trial (Hulse: 100, i.e.\ 1\,s at $\Delta t=0.01$\,s). \\
\texttt{alpha}         & $\alpha = 1/\tau_{\mathrm{corr}}$ & Mean-reversion rate. Hulse $\tau_{\mathrm{corr}}=0.12$\,s, so $\alpha\approx 8.33\;\mathrm{s}^{-1}$. \\
\texttt{sigma}         & $\sigma = \sigma_{\omega}\sqrt{2\alpha}$ & SDE diffusion. Hulse $\sigma_{\omega}=40^{\circ}/\mathrm{s}\;\Rightarrow\;\sigma\approx 163.3^{\circ}/\mathrm{s}^{3/2}$. \\
\texttt{sqrt\_dt}      & $\sqrt{\Delta t}$         & Hulse $\sqrt{0.01}=0.1$. \\
\texttt{sigma\_step}   & $\sigma\sqrt{\Delta t}$   & Per-step noise scale, $\approx 16.3^{\circ}$. \\
\texttt{omega}         & $\omega[b,t]$, shape $(B,T)$ & Velocity trajectory; pre-allocated to zeros. Hulse: $\omega[0]=0$. \\
\texttt{eta}           & $\eta[b,t]\sim\mathcal{N}(0,1)$, shape $(B,T)$ & The driving Gaussian noise, sampled once. \\
\bottomrule
\end{tabular}
\end{center}

\section*{Where these come from}

The two-step derivation is
\begin{equation}
\boxed{\;\;\sigma_{\omega}\;\xleftarrow{\text{user input}}\;
       \sigma=\sigma_{\omega}\sqrt{2\alpha}\;\xleftarrow{\text{calibrate}}\;
       \sigma\sqrt{\Delta t}\;\xleftarrow{\text{discretise}}\;
       \text{per-step Gaussian increment.}\;\;}
\end{equation}
The point of going through $\sigma$ rather than skipping straight to a
per-step std is that it makes the \emph{stationary} statistics
($\sigma_{\omega}$ = the std you'd measure on a long trajectory) match
Hulse's spec --- which is the meaningful biological quantity.

\section*{Numerical check}

Stationary std after one Euler step:
\begin{equation}
\mathrm{Var}\bigl(\omega_{\infty}\bigr)
   \;=\; \frac{(\sigma\sqrt{\Delta t})^{2}}{1-(1-\alpha\,\Delta t)^{2}}
   \;\approx\;
   \frac{\sigma^{2}\,\Delta t}{2\alpha\,\Delta t}
   \;=\; \frac{\sigma^{2}}{2\alpha}
   \;=\; \sigma_{\omega}^{2}.
\end{equation}
So $\mathrm{std}(\omega_{\infty})\;\approx\;\sigma_{\omega}=40^{\circ}/\mathrm{s}$,
matching Hulse's spec.

\section*{What happens next in the function}

The actual integration loop (immediately after the snippet) is the
discrete OU recursion
\begin{lstlisting}[language=Python]
decay = 1.0 - alpha * dt          # = 1 - dt/tau_corr  ~ 0.917 for Hulse
for t in range(1, T):
    omega[:, t] = decay * omega[:, t - 1] + sigma_step * eta[:, t]
\end{lstlisting}
which is exactly Eq.~(\ref{eq:ou-discrete}). After this loop, the
function inserts $\sim 20\%$ standing pauses (zeros into \texttt{omega}),
integrates \texttt{omega\_rad} to get true heading $\theta_{\mathrm{hd}}(t)$,
and packs the 3-channel input $u[t]=[\omega(t),\;\cos\theta_{0}\cdot
\mathbf{1}_{t=0},\;\sin\theta_{0}\cdot\mathbf{1}_{t=0}]$.

\section*{Figures (4 example trials, 4\,s each, $\Delta t=0.01$\,s)}

The three figures below are produced from
\texttt{generate\_path\_integration\_batch} itself --- not a re-implementation.

\begin{figure}[h]
\centering
\includegraphics[width=0.95\linewidth]{figures/ou_omega_theta_y.png}
\caption{\textbf{From OU velocity to network target.}
\emph{Top:} the OU angular velocity $\omega(t)$ in deg/s for four
trials. Grey shading marks a standing pause (only shown for trial~0):
$\omega$ is forced to zero, modelling the fly standing still.
\emph{Middle:} the unwrapped heading
$\theta_{\mathrm{hd}}(t)=\theta_0+\int\omega\,dt'$. During the stop, the
heading is held constant (flat plateau). Each trial starts at a random
$\theta_0\in[0,2\pi)$.
\emph{Bottom:} the 2-channel target $y(t)=(\cos\theta_{\mathrm{hd}},\sin\theta_{\mathrm{hd}})$
that the readout has to predict. Solid = $\cos$, dashed = $\sin$,
matched colour per trial.}
\end{figure}

\begin{figure}[h]
\centering
\includegraphics[width=0.95\linewidth]{figures/ou_u_channels.png}
\caption{\textbf{The 3-channel input $u[t]$ (trial 0).}
\emph{Top:} $u_0(t)=\omega(t)$ in deg/s --- the only time-varying
channel. Zeroed during the standing pause (grey).
\emph{Middle and bottom:} $u_1=\cos\theta_0$ and $u_2=\sin\theta_0$ are
\emph{single-frame impulses} at $t=0$ only (the orange/green spike), and
zero for all $t>0$. These encode the initial heading so the network can
anchor its bump at the correct phase before integrating velocity.}
\end{figure}

\begin{figure}[h]
\centering
\includegraphics[width=0.95\linewidth]{figures/ou_eta_and_hist.png}
\caption{\textbf{Driving noise $\eta$ and the resulting stationary
distribution of $\omega$.}
\emph{Left:} four independent draws of $\eta[t]\sim\mathcal{N}(0,1)$
(unit-variance white noise). This is what the OU update consumes via
$\sigma_{\mathrm{step}}\cdot\eta[t]$.
\emph{Right:} histogram of $\omega(t)$ pooled over 200 trials $\times$
900 post-burnin steps. Red curve is the target stationary
$\mathcal{N}(0,\sigma_{\omega}^{2})$ with $\sigma_{\omega}=40^{\circ}$/s.
The empirical std (32.3) lands a bit below the theoretical 40 because
$\sim$20\% of samples are zeroed by standing pauses --- exactly the
effect we want.}
\end{figure}

\end{document}
